\section{初等不定积分：从识别到算法}\label{ic:indefinite_classical_complete}

\begin{examtip}以下按被积函数的结构展开计算步骤；同类方法优先用最短换元，不把万能代换当作每题的第一步。\end{examtip}

\subsection{先整理：除法、配方、拆项与逆链式}\label{ic:indefinite_classical_complete:normalize}\begin{method}{先整理：除法、配方、拆项与逆链式}\textbf{识别信号} 分式分子次数不低于分母；分子接近分母导数。\par \textbf{操作步骤} 先长除；二次式配方；把分子拆成$A Q'(x)+B$；最后确认是否为$\frac{f'(x)}{f(x)}$或$f'(x)f(x)^r$。\par \textbf{疑难点与条件} 线性只适用于和及常数倍；积分不能分配到商或乘积。\end{method}

\begin{example}{考研风格｜长除与凑微分}{难度：$\star\star$}\label{ic:ex:indefinite_classical_complete:division}\textbf{题目} 求$\int\frac{x^3+2x^2+1}{x^2+1}\dd x$。\par \textbf{思路} 先除掉高次项，再拆导数项。\par \textbf{详细解答} 长除得$x+2+\frac{-x-1}{x^2+1}$，故原函数$\frac{x^2}{2}+2x-\frac12\ln(x^2+1)-\arctan x+C$。另一法先把分子写成$(x+2)(x^2+1)-(x+1)$，再逐项积分。\par \textbf{求导验算} 求导为$x+2-\frac{x+1}{x^2+1}$，通分恢复原分子。\par \textbf{疑难点与方法比较} 先作代数整理通常比立即换元省力。\end{example}

\begin{example}{考研风格｜导数配分母}{难度：$\star\star$}\label{ic:ex:indefinite_classical_complete:affine}\textbf{题目} 求$\int\frac{3x+4}{x^2+2x+5}\dd x$。\par \textbf{思路} $3x+4=\frac32(2x+2)+1$。\par \textbf{详细解答} 配方$x^2+2x+5=(x+1)^2+4$，答案$\frac32\ln(x^2+2x+5)+\frac12\arctan(\frac{x+1}{2})+C$。\par \textbf{求导验算} 导数$\frac{3x+3}{Q}+\frac1Q=\frac{3x+4}{Q}$。\par \textbf{疑难点与方法比较} 对数项与反正切项的系数分别由配导数和配方决定。\end{example}

\subsection{对数、反函数与整体换元}\label{ic:indefinite_classical_complete:logsub}\begin{method}{对数、反函数与整体换元}\textbf{识别信号} 出现$\frac{\dd x}{x}$配$\ln x$，或$\frac{\dd x}{1+x^2}$配$\arctan x$。\par \textbf{操作步骤} 将内层整体命名为$t$；同时替换自变量因子与微分；若换元后仍含$x$，用反函数回写。\par \textbf{疑难点与条件} 换元是局部可逆变换；不跨越对数、反三角或原分母的奇点。\end{method}

\begin{example}{考研风格｜反三角幂整体换元}{难度：$\star\star$}\label{ic:ex:indefinite_classical_complete:atanpower}\textbf{题目} 求$\int\frac{(\arctan x)^4\dd x}{1+x^2}$。\par \textbf{思路} 内层导数完整出现。\par \textbf{详细解答} 令$t=\arctan x$，$\dd t=\frac{\dd x}{1+x^2}$，积分$\frac{t^5}{5}+C=\frac{(\arctan x)^5}{5}+C$。\par \textbf{求导验算} 链式求导立即还原被积函数。\par \textbf{疑难点与方法比较} 反三角函数不总要分部；出现其导数时先换元。\end{example}

\begin{example}{考研风格｜对数换元后再分部}{难度：$\star\star\star$}\label{ic:ex:indefinite_classical_complete:xlog}\textbf{题目} 求$\int(\ln x)^2\dd x$，$x>0$。\par \textbf{思路} 可直接分部，也可令$t=\ln x$化指数多项式。\par \textbf{详细解答} 分部两次得$x(\ln x)^2-2\int\ln x\dd x=x[(\ln x)^2-2\ln x+2]+C$。另一法$x=\e^t$，积分$\int t^2\e^t\dd t=\e^t(t^2-2t+2)+C$，回代相同。\par \textbf{求导验算} 导数为$(\ln x)^2-2\ln x+2+2\ln x-2=(\ln x)^2$。\par \textbf{疑难点与方法比较} 换元可以把反函数积分转为指数多项式族。\end{example}

\subsection{分部积分：降阶、升阶与反向选择}\label{ic:indefinite_classical_complete:parts}\begin{method}{分部积分：降阶、升阶与反向选择}\textbf{识别信号} 乘积中一个因子求导会简化；或另一个因子积分会引入可用结构。\par \textbf{操作步骤} 先预测剩余积分，再选择$u,\dd v$；多项式反复求导；反三角或对数常求导；若遇循环，将未知积分移到同一侧。\par \textbf{疑难点与条件} 经验顺序不是定理，必须比较剩余积分是否变简单。\end{method}

\begin{example}{考研风格｜反向分部选择}{难度：$\star\star\star$}\label{ic:ex:indefinite_classical_complete:logsec}\textbf{题目} 求$\int x\ln(1+x^2)\dd x$。\par \textbf{思路} 令$t=1+x^2$再分部，比把$x$作为求导因子更短。\par \textbf{详细解答} $t=1+x^2$给$\frac12\int\ln t\dd t=\frac12[(1+x^2)\ln(1+x^2)-(1+x^2)]+C$。直接分部取$u=\ln(1+x^2),\dd v=x\dd x$也得$\frac{x^2}2\ln(1+x^2)-\int \frac{x^3}{1+x^2}\dd x$，长除后等价。\par \textbf{求导验算} 对第一式求导为$x\ln(1+x^2)$。\par \textbf{疑难点与方法比较} 先换元再分部可避免高次有理分式。\end{example}

\begin{example}{考研风格｜反三角函数分部}{难度：$\star\star\star$}\label{ic:ex:indefinite_classical_complete:arcsin}\textbf{题目} 求$\int\arcsin x\dd x$，$|x|<1$。\par \textbf{思路} 把隐含因子1积分成$x$。\par \textbf{详细解答} 取$u=\arcsin x,\dd v=\dd x$，积分$x\arcsin x-\int \frac{x}{\sqrt{1-x^2}}\dd x=x\arcsin x+\sqrt{1-x^2}+C$。第二法$x=\sin t$，对$\int t\cos t\dd t$分部，也得同式。\par \textbf{求导验算} 两项中的$\frac{x}{\sqrt{1-x^2}}$相消。\par \textbf{疑难点与方法比较} 最后根式前为正号；区间端点的可导性单独讨论。\end{example}

\subsection{循环分部与有限维待定原函数}\label{ic:indefinite_classical_complete:cyclic}\begin{method}{循环分部与有限维待定原函数}\textbf{识别信号} 指数乘三角、双曲，或求导后封闭于有限函数族。\par \textbf{操作步骤} 写联立积分关系，或设$F=\sum A_jf_j$并比较导数系数；解线性方程；最后检验共振或参数退化。\par \textbf{疑难点与条件} 循环回到原积分时必须移项，不能将中间积分误当结果。\end{method}

\begin{example}{考研风格｜循环分部与复指数三解}{难度：$\star\star\star$}\label{ic:ex:indefinite_classical_complete:expcycle}\textbf{题目} 求$\int\e^{2x}\sin3x\dd x$。\par \textbf{思路} 二次循环、待定系数或复指数。\par \textbf{详细解答} 设$I=\int\e^{2x}\sin3x$、$J=\int\e^{2x}\cos3x$。分部给$I=\frac{\e^{2x}\sin3x}{2}-\frac{3J}{2}$、$J=\frac{\e^{2x}\cos3x}{2}+\frac{3I}{2}$，解得$I=\frac{\e^{2x}(2\sin3x-3\cos3x)}{13}+C$。设$\e^{2x}(A\sin3x+B\cos3x)$比较也得$A=\frac{2}{13},B=-\frac{3}{13}$；复指数法取$\frac{\e^{(2+3i)x}}{2+3i}$虚部。\par \textbf{求导验算} 求导后正弦系数$\frac{4+9}{13}=1$，余弦系数抵消。\par \textbf{疑难点与方法比较} 待定原函数不是猜最终系数，而是解确定的线性方程组。\end{example}

\subsection{部分分式：重根、二次因子与不完全分解}\label{ic:indefinite_classical_complete:partial}\begin{method}{部分分式：重根、二次因子与不完全分解}\textbf{识别信号} 实有理式先化真分式。\par \textbf{操作步骤} 对$(x-a)^m$列全幂；对不可约二次$Q^m$每幂配置一次分子；可先分离特定因子，避免不必要的全分解。\par \textbf{疑难点与条件} 分解恒等式只在原函数定义的区间使用，不能跨极点合并常数。\end{method}

\begin{example}{考研风格｜一次重根的系统配置}{难度：$\star\star$}\label{ic:ex:indefinite_classical_complete:repeated}\textbf{题目} 求$\int\frac{\dd x}{x^2(x+1)}$。\par \textbf{思路} 重根$x^2$须保留$\frac{A}{x}+\frac{B}{x^2}$。\par \textbf{详细解答} $\frac{1}{x^2(x+1)}=-\frac{1}{x}+\frac{1}{x^2}+\frac{1}{x+1}$，故答案$-\ln|x|-\frac{1}{x}+\ln|x+1|+C$。\par \textbf{求导验算} 通分后的分子$-x(x+1)+(x+1)+x^2=1$。\par \textbf{疑难点与方法比较} 漏掉$\frac{1}{x}$项会使系数方程无解。\end{example}

\begin{example}{竞赛风格｜不可约二次重因子}{难度：$\star\star\star$}\label{ic:ex:indefinite_classical_complete:quadraticrepeated}\textbf{题目} 求$\int \frac{x^3\dd x}{(1+x^2)^2}$。\par \textbf{思路} 可用$t=1+x^2$，也可在原变量约化。\par \textbf{详细解答} $t=1+x^2$给$\frac12\int(t-1)t^{-2}\dd t=\frac12\ln(1+x^2)+\frac1{2(1+x^2)}+C$。另一法$\frac{x^3}{(1+x^2)^2}=\frac{x}{1+x^2}-\frac{x}{(1+x^2)^2}$，分别凑微分。\par \textbf{求导验算} 导数$\frac{x}{1+x^2}-\frac{x}{(1+x^2)^2}=\frac{x^3}{(1+x^2)^2}$。\par \textbf{疑难点与方法比较} 二次重因子遇奇次分子通常优先平方变量换元。\end{example}

\subsection{三角有理式的奇偶捷径}\label{ic:indefinite_classical_complete:trigparity}\begin{method}{三角有理式的奇偶捷径}\textbf{识别信号} 被积函数为$R(\sin x,\cos x)$。\par \textbf{操作步骤} 若对$\sin x$为奇函数，取$t=\cos x$；对$\cos x$为奇函数，取$t=\sin x$；同时变号不变，取$t=\tan x$；其余用$t=\tan(\frac{x}{2})$。\par \textbf{疑难点与条件} “奇”指整个有理表达式的符号变换性质；代换后用$\sin^2+\cos^2=1$化有理。\end{method}

\begin{example}{考研风格｜对正弦为奇的有理式}{难度：$\star\star$}\label{ic:ex:indefinite_classical_complete:sinodd}\textbf{题目} 求$\int\frac{\sin x\dd x}{1+\cos^2x}$。\par \textbf{思路} $t=\cos x$直接消去正弦微分。\par \textbf{详细解答} 积分$=-\int\frac{\dd t}{1+t^2}=-\arctan(\cos x)+C$。半角法也可有理化，但分母次数更高。\par \textbf{求导验算} 导数为$\frac{\sin x}{1+\cos^2x}$。\par \textbf{疑难点与方法比较} 先用符号对称性选变量，避免无条件使用万能代换。\end{example}

\begin{example}{考研风格｜同时变号不变}{难度：$\star\star\star$}\label{ic:ex:indefinite_classical_complete:simultaneous}\textbf{题目} 求$\int\frac{\dd x}{2\sin^2x+3\cos^2x}$。\par \textbf{思路} 除以$\cos^2x$后取$t=\tan x$。\par \textbf{详细解答} 在$\cos x\ne0$的区间，积分$=\int\frac{\dd t}{2t^2+3}=\frac1{\sqrt6}\arctan(\sqrt{\frac{2}{3}}\tan x)+C$。如需覆盖$\cos x=0$处，应在相邻区间调整常数，使原函数连续衔接。\par \textbf{求导验算} 导数为$\frac{\sec^2x}{3+2\tan^2x}$，乘$\cos^2x$恢复原式。\par \textbf{疑难点与方法比较} 原被积函数全实轴光滑，代换的极点不是原函数必有的奇点。\end{example}

\subsection{三角幂：留一个微分、降幂与递推}\label{ic:indefinite_classical_complete:trigpowers}\begin{method}{三角幂：留一个微分、降幂与递推}\textbf{识别信号} $\sin^mx\cos^nx$，或$\tan,\sec,\cot,\csc$的整数幂。\par \textbf{操作步骤} 奇幂留一个$\sin$或$\cos$；全偶用倍角；$\sec$偶幂留$\sec^2$，$\tan$奇幂配$\sec\tan$；无直接微分时分部递推。\par \textbf{疑难点与条件} 负幂也能递推，但分母零点必须从定义域剔除。\end{method}

\begin{example}{考研风格｜正切正割混合幂双解}{难度：$\star\star\star$}\label{ic:ex:indefinite_classical_complete:tansecpower}\textbf{题目} 求$\int\tan^3x\sec^4x\dd x$。\par \textbf{思路} 可留$\sec^2x$用$\tan x$，或留$\sec x\tan x$用$\sec x$。\par \textbf{详细解答} 取$t=\tan x$，得到$\int t^3(1+t^2)\dd t=\frac{t^4}{4}+\frac{t^6}{6}+C$。另一法$u=\sec x$，$\tan^2x=u^2-1$，积分$\int(u^2-1)u^3\dd u=\frac{u^6}{6}-\frac{u^4}{4}+C$，两式相差常数。\par \textbf{求导验算} 第一式求导为$(\tan^3+\tan^5)\sec^2=\tan^3\sec^4$。\par \textbf{疑难点与方法比较} 两条换元均可行时比较多项式次数。\end{example}

\begin{example}{考研风格｜不定积分降二阶递推}{难度：$\star\star\star$}\label{ic:ex:indefinite_classical_complete:sinrec}\textbf{题目} 推导$S_n=\int\sin^nx\dd x$并求$S_5$。\par \textbf{思路} 对$\sin^{n-1}x\sin x$分部。\par \textbf{详细解答} $S_n=-\sin^{n-1}x\cos x+(n-1)\int\sin^{n-2}x\cos^2x\dd x$，故$S_n=-\frac{\sin^{n-1}x\cos x}{n}+\frac{(n-1)S_{n-2}}{n}$。初值$S_0=x,S_1=-\cos x$。也可留$\sin x$令$u=\cos x$，直接得$S_5=-\cos x+\frac23\cos^3x-\frac15\cos^5x+C$。\par \textbf{求导验算} 多项式式求导为$\sin x(1-2\cos^2x+\cos^4x)=\sin^5x$。\par \textbf{疑难点与方法比较} 定积分Wallis递推是边界项消失后的特例，不定积分不能丢边界项。\end{example}

\subsection{双曲有理式与幂递推}\label{ic:indefinite_classical_complete:hyperbolic}\begin{method}{双曲有理式与幂递推}\textbf{识别信号} 由$\sinh,\cosh,\tanh$构成的有理式。\par \textbf{操作步骤} 通法$t=\e^x$；半角$t=\tanh(\frac{x}{2})$也可；幂函数用$\cosh^2-\sinh^2=1$，注意平方恒等式符号。\par \textbf{疑难点与条件} 实$\cosh x>0$，但$\sinh x=0$会产生分母奇点；$\sqrt{\sinh^2x}=|\sinh x|$。\end{method}

\begin{example}{考研风格｜双曲半角有理化}{难度：$\star\star\star$}\label{ic:ex:indefinite_classical_complete:hyperbolicrat}\textbf{题目} 求$\int\frac{\dd x}{2+\cosh x}$。\par \textbf{思路} 取$t=\tanh(\frac{x}{2})$，$|t|<1$。\par \textbf{详细解答} $\cosh x=\frac{1+t^2}{1-t^2}$、$\dd x=\frac{2\dd t}{1-t^2}$，积分$=2\int\frac{\dd t}{3-t^2}=\frac2{\sqrt3}\operatorname{artanh}(\frac{t}{\sqrt3})+C$。也可$t=\e^x>0$，积分化为$2\int\frac{\dd t}{t^2+4t+1}$，配方给对数形式。\par \textbf{求导验算} $\frac{\dd}{\dd x}\operatorname{artanh}(\frac{t}{\sqrt3})=\frac{t'}{\sqrt3(1-\frac{t^2}{3})}$，代$t'=\frac{1-t^2}{2}$还原原式。\par \textbf{疑难点与方法比较} 双曲型二次分母可能导出反双曲函数或实对数，不能误写成反正切。\end{example}

\begin{example}{竞赛风格｜双曲幂递推的符号}{难度：$\star\star\star$}\label{ic:ex:indefinite_classical_complete:hyperbolicpower}\textbf{题目} 令$C_n=\int\cosh^nx\dd x$，推导递推并求$C_3$。\par \textbf{思路} 分部并用$\sinh^2=\cosh^2-1$。\par \textbf{详细解答} $C_n=\cosh^{n-1}x\sinh x-(n-1)\int\cosh^{n-2}x\sinh^2x\dd x$，所以$C_n=\frac{\cosh^{n-1}x\sinh x}{n}+\frac{(n-1)C_{n-2}}{n}$。$C_3=\sinh x+\frac{\sinh^3x}{3}+C$，也由$t=\sinh x$得到。\par \textbf{求导验算} 导数$\cosh x(1+\sinh^2x)=\cosh^3x$。\par \textbf{疑难点与方法比较} 与正弦递推形式相近，但边界项符号不同。\end{example}

\subsection{Euler三替换的反解公式}\label{ic:indefinite_classical_complete:eulerthree}\begin{method}{Euler三替换的反解公式}\textbf{识别信号} $Q=ax^2+bx+c$为实二次式，$\sqrt Q$所在支路可微。\par \textbf{操作步骤} 若$a>0$设$\sqrt Q=\sqrt a\,x+t$，得$x=\frac{t^2-c}{b-2\sqrt a\,t}$；若$c>0$设$\sqrt Q=tx+\sqrt c$，得$x=\frac{2\sqrt c\,t-b}{a-t^2}$；若$Q=a(x-r)(x-s)$有实根，设$\sqrt Q=t(x-r)$，得$x=\frac{t^2r-as}{t^2-a}$。\par \textbf{疑难点与条件} 反解分母为0的参数需单独检查；平方后必须代回原式检验根式符号。三类方法至少一类可用于非空实正则区间。\end{method}

\begin{example}{考研风格｜第一Euler替换的完整计算}{难度：$\star\star\star$}\label{ic:ex:indefinite_classical_complete:eulerfirst}\textbf{题目} 求$\int\frac{\dd x}{\sqrt{x^2+2x+2}}$。\par \textbf{思路} 第一类替换，或先平移配方。\par \textbf{详细解答} 设$t=\sqrt{x^2+2x+2}-x$，则$x=\frac{t^2-2}{2-2t}$。由$\sqrt Q=x+t$微分，$\frac{(x+1)\dd x}{\sqrt Q}=\dd x+\dd t$，结合$x(t)$可化为$\frac{\dd x}{\sqrt Q}=-\frac{\dd t}{t-1}$，所以答案$-\ln|t-1|+C$。更短法$u=x+1$，结果$\ln(u+\sqrt{1+u^2})+C$。因$(\sqrt Q-x-1)(\sqrt Q+x+1)=1$，两式相同。\par \textbf{求导验算} $\frac{\dd\ln(x+1+\sqrt Q)}{\dd x}=\frac{1}{\sqrt Q}$。\par \textbf{疑难点与方法比较} 通法解释为何可有理化，配方通常是本题的优先方法。\end{example}

\begin{example}{考研风格｜第三Euler替换与比值根式}{难度：$\star\star\star$}\label{ic:ex:indefinite_classical_complete:eulerthird}\textbf{题目} 求$\int\frac{\dd x}{\sqrt{(x-1)(x+2)}}$，$x>1$。\par \textbf{思路} 以实根1为基点，设$t=\frac{\sqrt Q}{x-1}$。\par \textbf{详细解答} $t^2=\frac{x+2}{x-1}$，得$x=\frac{t^2+2}{t^2-1}$、$\dd x=-\frac{6t\dd t}{(t^2-1)^2}$、$\sqrt Q=\frac{3t}{t^2-1}$（$t>1$），故积分$=-2\int\frac{\dd t}{t^2-1}=\ln|\frac{t+1}{t-1}|+C$。也可先配方$Q=(x+\frac{1}{2})^2-(\frac{3}{2})^2$，得到$\ln(x+\frac{1}{2}+\sqrt Q)+C$。\par \textbf{求导验算} 对配方结果求导为$\frac{1}{\sqrt Q}$；与参数式只差常数。\par \textbf{疑难点与方法比较} 在$x<-2$支路须重新选择符号，不能照抄$t>1$。\end{example}

\subsection{二次根式正负幂的递推}\label{ic:indefinite_classical_complete:powerrec}\begin{method}{二次根式正负幂的递推}\textbf{识别信号} $\int(x^2+a^2)^{\frac{m}{2}}\dd x$，指数可逐步降到基本积分。\par \textbf{操作步骤} 对$x(x^2+a^2)^{\frac{m}{2}}$求导并重组；正整数与负整数分别从低阶初值递推；$m=-1$等退化位置改用基本对数式。\par \textbf{疑难点与条件} 递推公式的分母为0时不是积分不存在，而是该公式不适用。\end{method}

\begin{theorem}{统一幂递推}令$J_p=\int(x^2+a^2)^p\dd x$，则$p\ne-\frac{1}{2}$时$J_p=\frac{x(x^2+a^2)^p}{2p+1}+\frac{2pa^2}{2p+1}J_{p-1}$。\end{theorem}\begin{proof}$[x(x^2+a^2)^p]'=(2p+1)(x^2+a^2)^p-2pa^2(x^2+a^2)^{p-1}$，移项积分。$p=-\frac{1}{2}$的初值为$\ln(x+\sqrt{x^2+a^2})$（$a>0$）。\end{proof}

\begin{example}{竞赛风格｜根式三次幂的两解}{难度：$\star\star\star$}\label{ic:ex:indefinite_classical_complete:positivehalf}\textbf{题目} 求$\int(1+x^2)^{\frac{3}{2}}\dd x$。\par \textbf{思路} 统一递推或双曲换元。\par \textbf{详细解答} 递推$p=\frac{3}{2}$给$J_{\frac{3}{2}}=\frac{x(1+x^2)^{\frac{3}{2}}}{4}+\frac{3J_{\frac{1}{2}}}{4}$，而$J_{\frac{1}{2}}=\frac{x\sqrt{1+x^2}+\ln(x+\sqrt{1+x^2})}{2}$。故结果$\frac{x(2x^2+5)\sqrt{1+x^2}}8+\frac38\ln(x+\sqrt{1+x^2})+C$。另一法$x=\sinh t$，积分$\cosh^4t$后降幂。\par \textbf{求导验算} 对两项求导、通分后分子为$8(1+x^2)^2$，恢复$(1+x^2)^{\frac{3}{2}}$。\par \textbf{疑难点与方法比较} 正幂递推降低指数，负幂则需把恒等式反向求解。\end{example}
